\mxunit{automatismes-annexe-diag}{Annexe A — Diagnostic du 29 septembre 2026 et couverture par les fiches}
Données : test de positionnement 2G du 29/09/2026 (moyenne de classe en \%, nombre d'élèves sous $\pc{50}$ et à partir de $\pc{80}$). La colonne « Fiches » indique le nombre de fiches (sur $88$) où l'étiquette de diagnostic correspondante apparaît au moins une fois ; plusieurs attendus partagent une même étiquette.\par\smallskip
{\footnotesize\mxarraystretch{1.15}\begin{tabular}{@{}l l p{7.1cm} r r r l r@{}}\toprule
N° & Étiq. & Attendu & Moy. & $<\pc{50}$ & $\geqslant\pc{80}$ & Priorité & Fiches \\ \midrule
01 & CON & Comprendre un algorithme — programme de construction & \num{60,7} & 11 & 17 & À entretenir & 6 \\
02 & ALG & Comprendre un algorithme — programme de calcul & \num{42,9} & 16 & 12 & \textbf{Renf.\ prioritaire} & 17 \\
03 & NUM & Diverses représentations d'un même nombre (décimale, fractionnaire, notation scientifique, droite graduée) & \num{74,5} & 3 & 14 & À entretenir & 13 \\
04 & CAL & Calculer avec relatifs, fractions, décimaux & \num{65,4} & 6 & 8 & À entretenir & 21 \\
05 & PUI & Calculs simples avec puissances, notation scientifique & \num{67,1} & 7 & 15 & À entretenir & 6 \\
06 & SUB & Substituer dans une expression algébrique & \num{57,1} & 4 & 2 & À consolider & 8 \\
07 & DEV & Développer, factoriser, réduire (cas très simples) & \num{74,1} & 2 & 11 & À entretenir & 8 \\
08 & EQ & Résoudre algébriquement une équation du premier degré & \num{38,0} & 21 & 2 & \textbf{Renf.\ prioritaire} & 29 \\
09 & LIT & Calcul littéral~: démontrer, valider/réfuter une conjecture, traduire par une expression & \num{34,8} & 15 & 0 & \textbf{Renf.\ prioritaire} & 13 \\
10 & PRO & Problèmes de proportionnalité (pourcentages, échelles, agrandissement-réduction) & \num{52,0} & 13 & 1 & À consolider & 18 \\
11 & QP & Calculer une quatrième proportionnelle & \num{64,3} & 10 & 18 & À entretenir & 5 \\
12 & GEO & Figures, configurations, transformations~: déterminer des grandeurs & \num{35,5} & 24 & 0 & \textbf{Renf.\ prioritaire} & 28 \\
13 & STA & Effectifs, fréquences\,; moyenne, médiane, étendue & \num{51,8} & 8 & 9 & À consolider & 15 \\
14 & LEC & Lire et interpréter données brutes, tableaux, diagrammes & \num{57,1} & 12 & 6 & À consolider & 14 \\
15 & IMG & Image ou antécédent à partir d'une représentation & \num{76,2} & 6 & 14 & À entretenir & 12 \\
16 & MOD & Problèmes modélisés par des fonctions & \num{78,6} & 6 & 22 & À entretenir & 13 \\
17 & DEM & Raisonner et s'initier à la démonstration (figures, configurations) & \num{34,2} & 23 & 1 & \textbf{Renf.\ prioritaire} & 17 \\
18 & REP & Repérage (droite graduée, plan, pavé, sphère) & \num{32,1} & 19 & 9 & \textbf{Renf.\ prioritaire} & 15 \\
19 & SOL & Reconnaître des solides & \num{34,5} & 20 & 1 & \textbf{Renf.\ prioritaire} & 10 \\
20 & PNP & Reconnaître proportionnalité / non-proportionnalité & \num{61,9} & 10 & 8 & À entretenir & 5 \\
21 & CMP & Calculs et comparaisons pour traiter des problèmes & \num{52,7} & 6 & 2 & À consolider & 13 \\
22 & DIV & Problèmes de divisibilité (engrenages, conjonction de phénomènes) & \num{42,9} & 10 & 6 & \textbf{Renf.\ prioritaire} & 11 \\
23 & PER & Représentations de solides (perspective, face/dessus, sections, patrons) & \num{67,9} & 9 & 19 & À entretenir & 5 \\
24 & PRB & Probabilités en situation d'équiprobabilité & \num{57,1} & 12 & 16 & À consolider & 17 \\
25 & REG & Passer d'un mode de représentation d'une fonction à un autre & \num{25,0} & 21 & 7 & \textbf{Renf.\ prioritaire} & 21 \\
\bottomrule\end{tabular}}\par\bigskip
\textbf{Fréquence des questions « Retour » et « Fil rouge » par groupe de priorité} (nombre de fiches contenant au moins une question de ce groupe en positions 4–5) :\par\smallskip
\begin{tabular}{@{}l r@{}}\toprule Groupe & Fiches sur 88 \\ \midrule
Renforcement prioritaire & 88 \\
À consolider & 36 \\
À entretenir & 52 \\
\bottomrule\end{tabular}\par